\documentclass[10pt,reqno]{amsart}

\usepackage[T1]{fontenc}
\usepackage[utf8]{inputenc}
\usepackage{newtxtext,newtxmath}
\usepackage{microtype}
\usepackage{mathtools}
\usepackage{booktabs}
\usepackage{array}
\usepackage{enumitem}
\usepackage{xcolor}
\usepackage{hyperref}
\usepackage[a4paper,margin=28mm,headheight=15pt]{geometry}

\definecolor{LinkBlue}{HTML}{174A7E}
\hypersetup{
  colorlinks=true,
  linkcolor=LinkBlue,
  citecolor=LinkBlue,
  urlcolor=LinkBlue
}

\setlength{\parskip}{2pt}
\setlist[itemize]{leftmargin=6mm,itemsep=1.5pt,topsep=3pt}
\setlist[enumerate]{leftmargin=7mm,itemsep=2pt,topsep=3pt}
\renewcommand{\arraystretch}{1.15}
\emergencystretch=2em

\newtheorem{theorem}{Theorem}[section]
\newtheorem{proposition}[theorem]{Proposition}
\newtheorem{lemma}[theorem]{Lemma}
\newtheorem{corollary}[theorem]{Corollary}
\theoremstyle{definition}
\newtheorem{definition}[theorem]{Definition}
\newtheorem{remark}[theorem]{Remark}

\title[Orthogonal Join Reduction and Volume Rationality for One-Angle Spherical Bar Cycles]
{Co-Normalizer Orthogonal Join Reduction and Rational Chern--Simons Volume on One-Angle Spherical Bar Cycles}
\author{Argus Open Research Project}
\date{September 5, 2026}
\subjclass[2020]{Primary 57K32; Secondary 58J28, 19E20}
\begin{document}
\begin{abstract}
A conjecture in low-dimensional topology asks whether the $\mathrm{SU}(2)$ Chern--Simons invariant of every flat connection on a closed $3$-manifold is rational modulo~$1$. For the one-angle spherical bar-cycle construction in the Bloch/scissors congruence group, we prove that the co-normalizer flatness constraint forces an orthogonal spherical join, and that vanishing of the Dehn invariant on the resulting cycle forces its normalized Chern--Simons volume to be rational.
\end{abstract}
\maketitle
\section{Main Rationality Reduction}
\begin{theorem}
In the one-angle spherical bar-cycle family, imposing the co-normalizer holonomy constraint decomposes every admissible tetrahedron cycle as an orthogonal spherical join; tensoring the Schlafli/Dehn invariant over $\mathbb{Q}$ shows that whenever the Dehn invariant vanishes, the dihedral angles are rational multiples of $\pi$ and the normalized Chern--Simons invariant lies in $\mathbb{Q}/\mathbb{Z}$.
\end{theorem}
\end{document}
